% Part:sets-functions-relations
% Chapter: sets
% Section: reduction

\documentclass[../../../include/open-logic-section]{subfiles}

\begin{document}

\olfileid{sfr}{siz}{red}

\olsection{হ্রাসকরণ}

\begin{editorial}
  এই অংশে হ্রাসকরণের সাহায্যে অতালিকায়নযোগ্যতা প্রমাণ করা হয়েছে; ফলগুলি \olref[nen]{sec}-এর ফলের সঙ্গে মেলে। \olref[nen-alt]{sec}-এর ফলের সঙ্গে মেলে এমন একটি বিকল্প, কিছুটা সংক্ষিপ্ততর সংস্করণ \olref[red-alt]{sec}-এ দেওয়া আছে।
\end{editorial}

কর্ণীকরণ যুক্তির সাহায্যে আমরা দেখিয়েছি যে $\Pow{\PosInt}$ !!{nonenumerable}। সকল অসীম $0$-$1$ অনুক্রমের সেট $\Bin^\omega$ যে !!{nonenumerable}, তার একটি প্রমাণ আমাদের আগে থেকেই ছিল। $\Pow{\PosInt}$ যে !!{nonenumerable}, তা প্রমাণের আরেকটি উপায় এখানে: দেখাও যে \emph{$\Pow{\PosInt}$ যদি !!{enumerable} হয়, তবে $\Bin^\omega$-ও !!{enumerable}}। যেহেতু আমরা জানি $\Bin^\omega$ !!{enumerable} নয়, তাই $\Pow{\PosInt}$-ও তা হতে পারে না। এক সমস্যাকে অন্য সমস্যায় \emph{হ্রাস করা} বলতে এটিই বোঝায়—এই ক্ষেত্রে আমরা $\Bin^\omega$-কে তালিকায়ন করার সমস্যাটিকে $\Pow{\PosInt}$-কে তালিকায়ন করার সমস্যায় হ্রাস করি। দ্বিতীয়টির একটি সমাধান—$\Pow{\PosInt}$-এর একটি তালিকায়ন—প্রথমটির একটি সমাধান—$\Bin^\omega$-এর একটি তালিকায়ন—দিত।

কোনো সেট~$B$-কে তালিকায়ন করার সমস্যাকে কীভাবে সেট~$A$-কে তালিকায়ন করার সমস্যায় হ্রাস করব?~$A$-এর একটি তালিকায়নকে~$B$-এর তালিকায়নে বদলানোর একটি উপায় দিই। এর সবচেয়ে সহজ উপায় হলো !!a{surjective} অপেক্ষক $f\colon A \to B$ সংজ্ঞায়িত করা। $x_1$, $x_2$, \dots{} যদি~$A$-এর তালিকায়ন হয়, তবে $f(x_1)$, $f(x_2)$, \dots{} হবে~$B$-এর তালিকায়ন। আমাদের ক্ষেত্রে আমরা একটি সমাপতিত অপেক্ষক $f\colon \Pow{\PosInt} \to \Bin^\omega$ খুঁজছি।

\begin{prob}
দেখাও যে কোনো !!{injective} অপেক্ষক $g\colon B \to A$ থাকলে এবং~$B$ !!{nonenumerable} হলে~$A$-ও তাই।~$A$-এর একটি তালিকায়নকে~$B$-এর তালিকায়নে বদলাতে কীভাবে~$g$ ব্যবহার করা যায় তা দেখিয়ে প্রমাণটি করো।
\end{prob}

\begin{proof}[হ্রাসকরণের সাহায্যে {\olref[nen]{thm:nonenum-pownat}}-এর প্রমাণ]
ধরি, $\Pow{\PosInt}$ !!{enumerable}, এবং তাই তার একটি তালিকায়ন $Z_{1}$, $Z_{2}$, $Z_{3}$, \dots{} আছে।

$f(Z)$-কে সেই অনুক্রম $s$ নিই, যার জন্য $s(n) = 1$ যদি এবং কেবল যদি $n \in Z$, এবং অন্যথায় $s(n) = 0$; এভাবে অপেক্ষক $f \colon \Pow{\PosInt} \to \Bin^\omega$ সংজ্ঞায়িত করি। % OLSIZ-004 TARGET-CORRECTION-BEGIN
\par\smallskip\noindent\textnormal{\small\textbf{উৎস-সংশোধন (OLSIZ-004):} হিমায়িত ইংরেজি উৎসে ফল অনুক্রমটির নাম $s_k$, কিন্তু $k$ কোথাও বাঁধা বা নির্ধারিত নয়। বাংলা পাঠে একই ফলকে সর্বত্র বাঁধা নাম $s$ দিয়ে লেখা হয়েছে; তার $n$-তম মানই সদস্যতার নির্দেশক।} % OLSIZ-004 TARGET-CORRECTION-END
এটি স্পষ্টতই একটি অপেক্ষক সংজ্ঞায়িত করে, কারণ $Z \subseteq \PosInt$ হলে যেকোনো $n \in \PosInt$ হয় $Z$-এর !!a{element}, নয়তো নয়। যেমন, ধনাত্মক জোড় সংখ্যার সেট $2\PosInt = \{2,
4, 6, \dots\}$-কে অনুক্রম $010101\dots$-এ, শূন্য সেটকে $0000\dots$-এ, এবং $\PosInt$-কে $1111\dots$-এ পাঠানো হয়।

এটি !!{surjective}-ও: $0$ ও $1$-এর প্রত্যেক অনুক্রম কোনো ধনাত্মক পূর্ণসংখ্যার সেটের সঙ্গে সঙ্গতিপূর্ণ, যথা অনুক্রমের যেসব স্থানে~$1$ আছে সেই স্থানগুলির সঙ্গে সংশ্লিষ্ট পূর্ণসংখ্যাগুলি যার সদস্য। আরও নির্দিষ্টভাবে, ধরি $s \in \Bin^\omega$। নিচের নিয়মে $Z \subseteq \PosInt$ সংজ্ঞায়িত করি:
\[
Z = \Setabs{n \in \PosInt}{s(n) = 1}
\]
তাহলে~$f$-এর সংজ্ঞা দেখে যাচাই করা যায় যে $f(Z) = s$।

এখন তালিকাটি বিবেচনা করো:
\[
f(Z_1), f(Z_2), f(Z_3), \dots
\]
যেহেতু $f$ !!{surjective}, তাই $\Bin^\omega$-এর প্রত্যেক সদস্যকে কোনো আর্গুমেন্টের জন্য~$f$-এর মান হিসেবে আসতে হবে; ফলে সেটি তালিকাতেও আসবে। তাই এই তালিকা~$\Bin^\omega$-এর সবকিছুকে তালিকায়ন করে।

সুতরাং $\Pow{\PosInt}$ !!{enumerable} হলে $\Bin^\omega$ !!{enumerable} হতো। কিন্তু $\Bin^\omega$ !!{nonenumerable}
(\olref[nen]{thm:nonenum-bin-omega})। তাই $\Pow{\PosInt}$ !!{nonenumerable}।
\end{proof}

\begin{explain}
হ্রাস কোন অভিমুখে যায়, তা নিয়ে সহজেই বিভ্রান্তি হতে পারে। যেমন, !!a{surjective} অপেক্ষক $g \colon \Bin^\omega \to B$, $B$ যে !!{nonenumerable} তা \emph{প্রতিষ্ঠা করে না}। ($g(s) = s(1)$ দ্বারা নির্ধারিত $g
\colon \Bin^\omega \to \Bin$ বিবেচনা করো; অপেক্ষকটি $0$ ও $1$-এর অনুক্রমকে তার প্রথম !!{element}-এ পাঠায়। এটি !!{surjective}, কারণ কিছু অনুক্রম $0$ দিয়ে এবং কিছু অনুক্রম $1$ দিয়ে শুরু হয়। কিন্তু $\Bin$ সসীম।) আরও লক্ষ করো, অপেক্ষক~$f$-কে !!{surjective} হতেই হবে, নইলে যুক্তিটি কাজ করে না: $f(x_1)$, $f(x_2)$, \dots{}-এ~$B$-এর সব !!{element}s অবশ্যই থাকবে—এ নিশ্চয়তা তখন পাওয়া যেত না। যেমন,
\[
h(n) = \underbrace{000\dots0}_{\text{$n$টি $0$}}111\dots
\] % OLSIZ-005 TARGET-CORRECTION-BEGIN
\par\smallskip\noindent\textnormal{\small\textbf{উৎস-সংশোধন (OLSIZ-005):} হিমায়িত ইংরেজি উৎসের প্রদর্শিত ফলটি শুধু $n$টি শূন্যের একটি সসীম প্রতীকক্রম, তাই সেটি $\Bin^\omega$-এর সদস্য নয়। বাংলা পাঠে তার পরে অসীমসংখ্যক $1$ যোগ করা হয়েছে; এতে প্রতিটি $h(n)$ সত্যিই অসীম বাইনারি অনুক্রম হয়, অথচ যেমন সর্বশূন্য অনুক্রমটি এখনও $h$-এর মান নয়।} % OLSIZ-005 TARGET-CORRECTION-END
একটি অপেক্ষক $h\colon \PosInt \to
\Bin^\omega$ সংজ্ঞায়িত করে, কিন্তু $\PosInt$ !!{enumerable}।
\end{explain}

\begin{prob}
হ্রাসকরণ যুক্তির সাহায্যে দেখাও যে ধনাত্মক পূর্ণসংখ্যার ক্রমযুগলের সকল \emph{সেটের} সেট !!{nonenumerable}।
\end{prob}

\begin{prob}\label{sfr:siz:red:prob:nat-nat}
  হ্রাসকরণ যুক্তির সাহায্যে দেখাও যে সকল অপেক্ষক $f\colon \Nat \to \Nat$-এর সেট~$X$ !!{nonenumerable}। (ইঙ্গিত: $X$ থেকে~$\Bin^\omega$-তে একটি সমাপতিত অপেক্ষক দাও।)
\end{prob}

\begin{prob}
হ্রাসকরণ যুক্তির সাহায্যে দেখাও যে স্বাভাবিক সংখ্যার অসীম অনুক্রমগুলির সেট $\Nat^\omega$ !!{nonenumerable}।
\end{prob}

\begin{prob}
ধরি, $P$ হলো ধনাত্মক পূর্ণসংখ্যার সেট থেকে $\{0\}$ সেটে সকল অপেক্ষকের সেট, এবং $Q$ হলো ধনাত্মক পূর্ণসংখ্যার সেট থেকে $\{0\}$ সেটে সকল \emph{আংশিক} অপেক্ষকের সেট। দেখাও যে $P$ !!{enumerable}, কিন্তু $Q$ নয়। (ইঙ্গিত: $\Bin^\omega$-কে তালিকায়ন করার সমস্যাটিকে~$Q$-কে তালিকায়ন করার সমস্যায় হ্রাস করো।)
\end{prob}

\begin{prob}
ধরি, $S$ হলো ধনাত্মক পূর্ণসংখ্যার সেট থেকে \{0,1\} সেটে সকল !!{surjective} অপেক্ষকের সেট, অর্থাৎ $S$-এর সদস্য হল সকল !!{surjective}~$f\colon \PosInt \to \Bin$। দেখাও যে $S$ !!{nonenumerable}।
\end{prob}

\begin{prob}
দেখাও যে সকল বাস্তব সংখ্যার সেট~$\Real$ !!{nonenumerable}।
\end{prob}

\end{document}
